\section{Method}
\label{sec:method}

Given a synchronized stereo video stream
$\{I_L^t, I_R^t\}_{t=1}^{T}$ from a calibrated rig with intrinsics
$\Kmat_L, \Kmat_R$, relative pose $(\Rmat_{LR}, \tvec_{LR})$, and baseline
$b = \|\tvec_{LR}\|$, \StereoSplat produces a dynamic set of 3D Gaussians
$\mathcal{G}(t)$ in the \emph{metric} world frame that can be rendered at
any viewpoint and any timestamp. Figure~\ref{fig:pipeline} gives an
overview. We retain the encoder, probabilistic decoder, and bidirectional
deformation field of \StreamSplat~\cite{wu2026streamsplat} and replace
three components: the depth source (Sec.~\ref{sec:method-depth}), the
camera model (Sec.~\ref{sec:method-raster}), and the position
interpretation (Sec.~\ref{sec:method-head}). We then describe the
streaming display module (Sec.~\ref{sec:method-display}) and the training
losses (Sec.~\ref{sec:method-loss}).

\begin{figure*}[t]
\centering
\framebox[0.95\textwidth]{\parbox{0.92\textwidth}{\centering\vspace{2pt}
\textbf{Pipeline.} Left/right frames and stereo metric depth enter the
reused \StreamSplat encoder ($V{=}2$); the Gaussian decoder output is
lifted to metric world coordinates by the Metric Head; a perspective
rasterizer with per-frame poses renders both views; a streaming display
module applies motion-adaptive EMA and opacity-stability filtering before
the final frame is written. (Figure to be rendered from a trained
checkpoint in the camera-ready version.)\vspace{2pt}}}
\caption{\StereoSplat overview. Green blocks are reused from
\StreamSplat~\cite{wu2026streamsplat}; orange blocks are new.}
\label{fig:pipeline}
\end{figure*}

\paragraph{Preliminaries: \StreamSplat.}
\StreamSplat encodes each frame into a canonical orthographic Gaussian
set with a transformer over $8{\times}8$ patches, predicts per-pixel
Gaussian attributes (position, scale, rotation, opacity, color) with a
probabilistic truncated-Gaussian sampler, and models inter-frame motion
with a bidirectional polynomial deformation field. The renderer is the
orthographic diff-gaussian-rasterizer with a fixed identity camera.
Position predictions are bounded to $[-1,1]$ and interpreted in the
canonical space, which normalizes scale away.

\subsection{Metric Depth from Stereo}
\label{sec:method-depth}

We replace the Depth-Anything~\cite{yang2024depth} monocular prior with a
stereo-matching front end. For each frame pair $(I_L^t, I_R^t)$ we run
IGEV-Stereo~\cite{xu2023igev} to obtain a disparity map
$\disp^t$ and a confidence map $\rho^t$. Metric depth follows from the
pinhole triangulation formula:
\begin{equation}
\label{eq:triang}
\depth^t(u,v) \;=\; \frac{b \, \focal_x}{\disp^t(u,v)},
\end{equation}
where $\focal_x$ is the left-camera focal length in pixels. Pixels with
$\rho^t < \tau_\rho$ or $\disp^t < \epsilon$ are masked out of all
downstream depth supervision. The depth is stored as a
$\mathtt{float32}$ metric map and consumed both as an encoder input
channel (replacing the relative-depth channel of \StreamSplat) and as a
supervision target (Sec.~\ref{sec:method-loss}). Disparities can be
precomputed offline (\texttt{preprocess\_disparity\_stereo.py}) or
computed online; both paths feed the same provider interface.

\subsection{Perspective Rasterization with Per-Frame Poses}
\label{sec:method-raster}

We replace the orthographic rasterizer with a perspective one that
consumes the true intrinsics and a per-frame camera pose. Concretely, for
a camera with intrinsics $\Kmat$ and camera-to-world pose
$(\Rmat_{c2w}, \tvec_{c2w})$ we build
\begin{equation}
\Vmat = \texttt{World2View}(\Rmat_{c2w}^\top, -\Rmat_{c2w}^\top \tvec_{c2w}),\quad
\Pmat = \texttt{Perspective}(\Kmat, H, W),
\end{equation}
and pass $\Vmat \Pmat$ to the diff-gaussian-rasterizer CUDA kernel. The
dynamic-motion extensions of the original orthographic renderer
(\texttt{scales\_t}, \texttt{rotations\_r}, time-basis polynomial motion,
time-dependent opacity decay) are preserved, so the existing dynamic
decoder operates unchanged. For a static rig, the left pose is the
identity and the right pose is $(\Rmat_{LR}, \tvec_{LR})$; for a moving
rig, per-frame poses come from an offline SLAM/VIO trajectory
(Sec.~\ref{sec:method-head}).

\subsection{Metric Back-Projection Head}
\label{sec:method-head}

The decoder outputs per-pixel Gaussian attributes in the network's
normalized convention. We reinterpret the three position channels as
$(\delta u, \delta v, \delta z_\text{rel})$ --- a sub-pixel offset and a
multiplicative depth residual --- and lift each Gaussian into the metric
world frame:
\begin{align}
\label{eq:backproj}
Z &= \depth(u,v) \cdot (1 + \gamma\, \delta z_\text{rel}), \\
X &= (u + \delta u - c_x)\, Z / \focal_x, \quad
Y = (v + \delta v - c_y)\, Z / \focal_y, \\
\mean_\text{world} &= \Rmat_{c2w}\, (X, Y, Z)^\top + \tvec_{c2w}.
\end{align}
Gaussian scales are likewise lifted to meters: the in-plane scales are
mapped through a sigmoid to a fixed range $[s_\text{min}, s_\text{max}]$
in meters, and the depth-axis scale is set proportional to $Z$ so that
distant Gaussians are appropriately larger. Rotation and color are kept
as predicted. This head adds no learnable parameters; it is a fixed
geometric transform that converts the network's normalized output into a
metric representation.

\subsection{Streaming Display Module}
\label{sec:method-display}

Dynamic GS methods suffer from temporal flickering, popping, and
high-frequency motion blur (Sec.~\ref{sec:related}). We address these
with a streaming display module that wraps the rasterizer in a ring
buffer of recent Gaussian frames and applies two training-free operators
at each step.

\paragraph{Motion-Adaptive Temporal EMA.}
Let $\mean_t, \rgb_t, \opac_t$ be the metric means, colors, and opacities
of the current frame, and $\mean_{t-1}, \rgb_{t-1}, \opac_{t-1}$ those of
the previous frame. We compute a per-Gaussian motion magnitude
$m_t = \|\mean_t - \mean_{t-1}\|$ and a mixing factor
$\alpha = \alpha_\text{static}\, \mathbb{1}[m_t < \tau_m] +
\alpha_\text{dyn}\, (1 - \mathbb{1}[m_t < \tau_m])$ with
$\alpha_\text{static} \gg \alpha_\text{dyn}$. The smoothed attributes are
\begin{equation}
\bar{\mean}_t = \alpha \mean_{t-1} + (1-\alpha) \mean_t, \quad
\bar{\rgb}_t = \alpha \rgb_{t-1} + (1-\alpha) \rgb_t,
\end{equation}
and similarly for $\opac$. Scale and rotation are not smoothed. This
strongly stabilizes static regions (where flicker is most visible) while
preserving motion detail under fast motion.

\paragraph{Opacity-Stability Filter.}
To suppress popping, we clamp the per-frame opacity delta and apply a
soft EMA: $\hat{\opac}_t = \beta\,(\opac_{t-1} + \text{clip}(\opac_t -
\opac_{t-1}, -\Delta, \Delta)) + (1-\beta)\, \opac_t$, with
$\Delta = 0.3$ and $\beta = 0.5$ in all experiments. This prevents
Gaussians from flashing in or out within a single frame.

\subsection{Training Losses}
\label{sec:method-loss}

The total loss combines photometric, stereo-consistency, metric-depth,
and perceptual terms:
\begin{equation}
\label{eq:loss}
\mathcal{L} = \Lphoto + \lambda_\text{stereo} \Lstereo +
\lambda_\text{depth} \Ldepth + \lambda_\text{lpips} \Llpips,
\end{equation}
where $\Lphoto = \|\hat{I}_L - I_L\|_1$ is the left-view photometric L1,
$\Lstereo = \|\hat{I}_R - I_R\|_1$ is the right-view photometric L1
(geometrically grounding the second view),
$\Ldepth = \|\hat{\depth} - \depth\|_1 \cdot \mathbb{1}[\rho > \tau_\rho]$
is the metric-depth L1 over confident pixels, and $\Llpips$ is the
LPIPS~\cite{zhang2018lpips} perceptual loss activated after a warmup.
Training is two-stage as in \StreamSplat: Stage~1 trains the static
encoder, Stage~2 freezes it and trains the dynamic decoder. We reuse
\StreamSplat's probabilistic sampling and adaptive opacity fusion
unchanged.
